% Proposed A.85. Source only until independently reviewed and privately built.
% Existing A_involutions.tex and all checkpoint27 research sources are unchanged.
\statement{A.85}{Proved}{The actual canonical Fox defect is induced under the joint Klein-four action}
Let \(\alpha=(-1,f)\in\widehat{GT}_0\) be a nonidentity
involution, and let \(u\in F'\) be its unique canonical parameter,
where \(F=\widehat F(x,y)\). Thus
\[
 \alpha(x)=x^{-1},\quad \alpha(y)=f^{-1}y^{-1}f,
 \qquad u=f\alpha(u),\qquad \iota(u)=u^{-1},
 \tag{A.85.1}
\]
and \(\iota\) interchanges x and y. Put
\[
 \begin{gathered}
 w=u^{-1}yu,\qquad \psi(x)=x,\quad\psi(y)=w,
 \qquad H=\overline{\langle x,w\rangle},\\
 v=\operatorname{Inn}(u^{-1})\iota,\qquad
 \tau=\alpha v,\qquad V=\langle\alpha,v\rangle.
 \end{gathered}
 \tag{A.85.2}
\]
Then V is the actual Klein-four group acting on F and H.
For every prime p use the full coefficient ring
\(\Lambda=\mathbb Z_p[[F]]\), and define
\[
 \begin{gathered}
 B=1+(y-1)u_y,\qquad
 L_H=\Lambda(x-1)+\Lambda(w-1),\\
 C_p=I_F/L_H
     \simeq\ker\bigl(\mathbb Z_p[[F/H]]\longrightarrow\mathbb Z_p\bigr).
 \end{gathered}
 \tag{A.85.3}
\]
The derivatives are the continuous left Fox derivatives specified below.
The actual coset space \(X=F/H\) satisfies
\[
       X^\alpha=X^v=X^\tau=\{H\}.
 \tag{A.85.4}
\]
There is a closed pointed transversal \((Y,H)\) for its nontrivial
V-orbits, and a compact \(\mathbb Z_p\)-module
\(D_p=\ker(\mathbb Z_p[[Y]]\to\mathbb Z_p)\), such that
\[
       C_p\simeq\mathbb Z_p[V]\otimes_{\mathbb Z_p}D_p
       \quad\hbox{as compact \(\mathbb Z_p[V]\)-modules}.
 \tag{A.85.5}
\]
In particular its Tate cohomology vanishes for V and every subgroup
of V. The complete coupled full-ring Fox identities are
(A.85.8)--(A.85.11) below. For any one p,
\[
       B\in\Lambda^\times
       \quad\Longleftrightarrow\quad C_p=0
       \quad\Longleftrightarrow\quad H=F.
 \tag{A.85.6}
\]
Neither (A.85.5) nor the coupled identities prove these equivalent
conditions. The induced description is over the constant finite-group
ring; no \(\Lambda\)-linear splitting or induced description is asserted.

\begin{proof}
\emph{The actual action and the full coefficient scope.}\par
Canonical uniqueness gives (A.85.1), with the swap duality at the
established \(\widehat{GT}_0\) scope. Direct substitution gives
\[
 \begin{gathered}
 v^2=\alpha^2=1,\qquad \alpha v=v\alpha,
 \qquad v(x)=w,\quad v(w)=x,\quad \alpha(w)=w^{-1},\\
 \alpha\psi=\psi c,\qquad v\psi=\psi\iota,
 \qquad\tau\psi=\psi(c\iota).
 \end{gathered}
 \tag{A.85.7}
\]
For example \(\alpha(u)=f^{-1}u\), and this identity, together
with the two-cycle identity for f, gives the commuting actions on
both free generators. Their three nonidentity homology actions are
respectively minus the identity, interchange, and negative interchange;
hence V has order four. All maps retain the full F. Finite V-stable
heads are cofinal by intersection of the four translates of any open
normal subgroup. In particular no maximal pro-p replacement is made.

Use the conventions
\[
 \begin{gathered}
 g-1=g_x(x-1)+g_y(y-1),\qquad
 (gh)_i=g_i+g h_i,\\
 J_\phi=(\partial_j\phi(x_i))_{i,j},\qquad
 J_{\phi\circ\rho}=\phi(J_\rho)J_\phi.
 \end{gathered}
\]
The augmentation ideal is the free compact left \(\Lambda\)-module
on \(x-1,y-1\). Thus row coefficient vectors multiply these
Jacobians on the right. Define
\[
 \begin{aligned}
 a&=u_x,& b&=u_y,& A_f&=1+(y-1)f_y,\\
 P&=u^{-1}(y-1)a,&D&=u^{-1}B,\\
 E&=u^{-1}\bigl(1+(x-1)a\bigr),&Q&=u^{-1}(x-1)b.
 \end{aligned}
\]
The full derivatives give
\[
 \begin{gathered}
 J_\psi=\begin{pmatrix}1&0\\P&D\end{pmatrix},\qquad
 J_v=\begin{pmatrix}P&D\\E&Q\end{pmatrix},\\
 J_\alpha=
 \begin{pmatrix}
 -x^{-1}&0\\
 -f^{-1}y^{-1}(y-1)f_x&-f^{-1}y^{-1}A_f
 \end{pmatrix},
 \qquad J_\tau=\alpha(J_v)J_\alpha.
 \end{gathered}
 \tag{A.85.8}
\]
Differentiating the two canonical equations in (A.85.1) also gives
\[
 \begin{gathered}
 \iota(a)=-u^{-1}b,\qquad\iota(b)=-u^{-1}a,\\
 (a,b)=(f_x,f_y)+f\,\alpha((a,b))J_\alpha.
 \end{gathered}
 \tag{A.85.9}
\]
There is no inversion of an individual Fox entry in this calculation.
The entire semilinear equivariance system is
\[
 \begin{aligned}
 \alpha(J_\psi)J_\alpha
     &=\begin{pmatrix}-x^{-1}&0\\0&-w^{-1}\end{pmatrix}J_\psi,\\
 v(J_\psi)J_v
     &=\begin{pmatrix}0&1\\1&0\end{pmatrix}J_\psi,\\
 \tau(J_\psi)J_\tau
     &=\begin{pmatrix}0&-w^{-1}\\-x^{-1}&0\end{pmatrix}J_\psi,\\
 \alpha(J_v)J_\alpha&=v(J_\alpha)J_v.
 \end{aligned}
 \tag{A.85.10}
\]
The third line is the product of the first two with their actual
coefficient actions. These are not identities of matrices with
commuting entries. In particular the second line yields
\[
 \begin{gathered}
 v(P)P+v(D)E=1,\qquad v(P)D+v(D)Q=0,\\
 B=\theta(B)A_f,\qquad
 \theta=\operatorname{Inn}(yf)\alpha,
 \qquad \theta^2=1,\quad\theta(A_f)A_f=1.
 \end{gathered}
 \tag{A.85.11}
\]
For the last line use the lower-right entry of the first line in
(A.85.10); the identities \((yf)\alpha(yf)=1\) and
\(\alpha(f)=f^{-1}\) fix its coefficient order. The matrices
\(J_\alpha,J_v\) are units because the underlying maps are
automorphisms; triangular invertibility makes \(A_f\) a unit.
The displayed identities involving all four entries of \(J_v\)
do not make D or B a unit.

\emph{The actual quotient module.}\par
The free compact module with marked basis \(e_x,e_w\) has the
semilinear action
\[
 \begin{gathered}
 \alpha(e_x)=-x^{-1}e_x,\quad
 \alpha(e_w)=-w^{-1}e_w,\qquad
 v(e_x)=e_w,\quad v(e_w)=e_x.
 \end{gathered}
\]
The involution and commuting relations follow on this basis from
(A.85.7). Sending the basis to \(x-1,w-1\) is equivariant, with
matrix \(J_\psi\). Its image is \(L_H\). Identifying the
corresponding translated cosets gives exactly the compact coset
quotient in (A.85.3), so the cokernel has this actual V-action.

A proper closed H is contained in a proper open subgroup of F, whose
finite coset module detects a nonzero element of the reduced compact
module. Hence \(C_p=0\) exactly when \(H=F\). If \(C_p=0\),
the endomorphism represented by \(J_\psi\) is surjective on a
rank-two free compact \(\Lambda\)-module. Its reductions over
\((\mathbb Z/p^n)[F/U]\) are surjections of finite sets of equal
order, so are injective. These heads separate the module; compactness
then gives a continuous inverse. The triangular matrix in (A.85.8)
is invertible exactly when D, equivalently B, is a unit. The converse
is immediate. This proves (A.85.6) without a local-ring assumption.

\emph{Section classes in the actual canonical subgroup.}\par
H is a closed subgroup of the full free profinite F, hence projective.
Its two marked generators give an epimorphism
\(\widehat{\mathbb Z}^2\to H_{\mathrm{ab}}\) whose composite
with \(H_{\mathrm{ab}}\to F_{\mathrm{ab}}\) is the identity.
Thus \(H_{\mathrm{ab}}=\widehat{\mathbb Z}^2\) with the named
marking. The v- and \(\tau\)-actions on its mod-two cohomology
are each the regular two-dimensional \(\mathbb F_2[C_2]\)-module.

For \(s=v,\tau\), put \(E_{H,s}=H\rtimes\langle s\rangle\).
The ambient \(F\rtimes\langle\tau\rangle\) is the actual
regular pants extension of C.43. For v the map which is the identity
on F and sends v to \(u^{-1}\iota\) identifies its ambient
extension with
\[
       F\rtimes\langle\iota\rangle
          \simeq C_2(\iota)\amalg\widehat{\mathbb Z}(x).
\]
The full free-product torsion and centralizer theorem therefore makes
all its involutions self-centralizing, as it does in C.43's regular
extension; see \citet[Theorems A' and B']{A_HR1985}. The same
self-centralizers hold in the closed subgroups \(E_{H,s}\).
Their index-two subgroup is projective. The real-projectivity
criterion and the allowed-class argument used with full proofs in
A.82 apply to these groups as well
\citep[Theorem~2.1]{A_PalQuickFormal2025}
\citep[Definition~1.1(4), Remark~1.2(b), Proposition~2.3]{AHaran1990}.
In particular any finite list of distinct involution classes admits
degree-one mod-two characters with independently prescribed values.
One first refines the finite head to separate the listed compact
classes, and allows both values on every unlisted head; thus no
proper solution of an embedding problem is being assumed.

The split Lyndon--Hochschild--Serre sequence for
\(H\to E_{H,s}\to C_2\) has only rows \(q=0,1\).
All positive cohomology of its regular \(q=1\) module is zero.
The possible transgression to \(H^2(C_2,\mathbb F_2)\) is zero
because the extension has the given section. Consequently
\[
          \dim_{\mathbb F_2}H^2(E_{H,s},\mathbb F_2)=1.
 \tag{A.85.12}
\]
For a list of distinct involution classes, choose the characters
which select its individual members. Their squares restrict to
linearly independent degree-two classes on that list of order-two
subgroups. Equation (A.85.12) permits at most one class; the actual
s supplies one. It maps to the unique ambient class.

For \(s=\alpha\), the subgroup \(H\rtimes\langle\alpha\rangle\)
is generated by the three actual involutions
\(\alpha,x\alpha,w\alpha\). Thus its degree-one mod-two
cohomology has dimension at most three. The same prescribed-character
argument permits at most three involution classes. Their three
displayed images are distinct already in
\(F_{\mathrm{ab}}/2\rtimes C_2\); C.28 says that these are
all the ambient classes. Hence this class map is a bijection as well.

Now suppose \(gH\) is fixed by \(s=\alpha,v\), or \(\tau\).
Then \(g^{-1}s(g)\in H\), and
\((g^{-1}s(g))s\) is an involution in \(E_{H,s}\) whose
ambient class is that of s. The class comparison just proved makes
it H-conjugate to s. Adjusting g on the right by that H-conjugator
gives an element of \(C_F(s)\). Each of these centralizers is
trivial: use C.28, C.43, or the full free-product description above,
respectively. Hence \(gH=H\), proving (A.85.4). This is a full
compact-coset assertion, not a statement about fixed points in every
finite coset quotient.

\emph{A closed pointed transversal and the precise induced module.}\par
The finitely generated F has a countable cofinal collection of open
normal subgroups, so X is compact metrizable. Choose decreasing
V-invariant clopen neighborhoods \(U_n\) of H with
\(U_0=X\) and \(\bigcap_nU_n=\{H\}\). Each shell
\(U_n\setminus U_{n+1}\) is a compact zero-dimensional free
V-space by (A.85.4). It has a clopen transversal: choose disjoint
translates of a sufficiently small clopen neighborhood at every
point, take a finite cover of the quotient, refine it to a disjoint
clopen partition, and select one sheet over every part.

Choose such a transversal \(T_n\) in each shell and set
\(Y=\{H\}\cup\bigcup_nT_n\). Its complement is the union
of the other open sheets, so Y is closed. Multiplication induces
a pointed V-equivariant continuous bijection
\[
       (V\times Y)/(V\times\{H\})\longrightarrow X.
 \tag{A.85.13}
\]
It is a homeomorphism by compactness. Splitting the free compact
module of Y by its distinguished point, its universal property
gives (A.85.5): send \(g\otimes(y-H)\) to \(gy-H\).
The finite V factor makes the tensor a finite direct sum of compact
modules. It is induced and coinduced from the trivial subgroup;
the standard finite-group cochain contraction, or Shapiro in Tate
cohomology, proves the asserted vanishing. Restricting to a subgroup
of V gives a finite sum of the same induced form.

The selected transversal is not asserted to be F-stable. Thus this
construction is over \(\mathbb Z_p[V]\), and gives no
\(\Lambda\)-linear splitting. In particular it does not prove
\(D_p=0\), which by (A.85.6) would settle the actual marked
generation gate. All uses of swap duality remain at the actual
negative-\(\widehat{GT}_0\) scope.
\end{proof}
